\section{Purpose and principal statements}\label{sec:intro}

The common question is how a geometric Hamiltonian, a positive theta
source, and a signed Fourier observation fit together. The source and its
transform are
\begin{align}
 \Phi(u)&=\sum_{m\ge1}
 (2\pi^2m^4\ee^{9u}-3\pi m^2\ee^{5u})\ee^{-\pi m^2\ee^{4u}},\label{eq:source}\\
 H_\lambda(z)&=\int_0^\infty\ee^{\lambda u^2}\Phi(u)\cos(zu)\dd u.
 \label{eq:Hintro}
\end{align}
We use the normalization
$H_0(z)=\xi((1+\ii z)/2)/8$, where
$\xi(s)=\tfrac12s(s-1)\pi^{-s/2}\Gamma(s/2)\zeta(s)$.
These are the conventions of Rodgers--Tao
\cite[Introduction, (1)--(4)]{RodgersTao2020}.
The canonical $\Phi$ is smooth, strictly positive and even. Its rapid
decay permits every finite real $\lambda$, and
$\partial_\lambda H_\lambda=-\partial_z^2H_\lambda$.

The shared geometric model is
\begin{equation}\label{eq:modelintro}
 C=\R_v\times S^2_{1/2},\qquad
 \Dc=-\ii\Gamma\partial_v+2D,\qquad
 L_0=\Dc^2=-\partial_v^2+4D^2.
\end{equation}
Here $D$ is the unit-sphere Dirac operator and $\Gamma$ its chirality.
The radial kinetic action multiplies Fourier amplitudes by $k^2$;
Newman source multiplication by $u^2$ differentiates them by
$-\partial_k^2$. This distinction motivates the two constructions.
Neither identifies free kinetic time with Newman time.

For the direct construction, Theorem~\ref{thm:direct} gives
\begin{equation}\label{eq:directintro}
 H_{1/2-\gamma t}(k)=H_{1/2}(0)\,[q_+(t,k)-q_-(t,k)],
 \quad 0\le t\le\gamma^{-1},\quad k\in\R.
\end{equation}
The $q_\pm$ are joint probabilities for no count followed by one of two
bounded cosine effects. The preparation and access to the transported
quadrature are specified inputs. For the auxiliary construction,
Theorem~\ref{thm:aux} gives
\begin{equation}\label{eq:auxintro}
 \Obs_{\lambda,\tau}(z)
 =\ee^{-\tau z^2}\frac{Q(\pi+4\tau)}{Q(\pi)}H_\lambda(z),
 \qquad Q(\beta)=\sum_{n\ge1}\ee^{-\beta n^2},\quad\tau\ge0.
\end{equation}
Its source is a linear amplitude and its observation is a stipulated
phase-sensitive pairing. These are different observed objects; the
normalizations cannot be interchanged.

The shared comparison and both branch constructions have written proofs
and internal acceptance within their stated scopes
\cite{TNDParent,TNDSharedReview,TNDDirect,TNDDirectReview,TNDAux,TNDAuxReview}.
This paper is a background synthesis of those results. It does not select
the next estimating problem for either branch. Section~\ref{sec:scope}
records the precise outstanding questions and the status as of 6 November
2026. No numerical evidence is used in the proofs below.

\begin{table}[ht]
\centering\small
\begin{tabularx}{\textwidth}{@{}lY@{}}\toprule
Symbol & Meaning \\\midrule
$u$; $v$, $X$ & Source coordinate; physical radial position and its multiplication operator. \\
$\beta$; $s$ & Angular damping; optional excited-shell filter strength. \\
$\lambda$ & Newman source-reweighting parameter. \\
$t$; $\tau$ & Laboratory time in the direct model; nonnegative heat/proper time in the auxiliary model. \\
$k$; $z$ & Real radial Fourier label; complex continuation of the observation parameter. \\
$m$; $n$; $\ell$ & Source-sum index; physical spinor shell; scalar harmonic degree. \\
$Z_\lambda$ & Full-line source mass $\int_\R\ee^{\lambda u^2}\Phi(u)\dd u=2H_\lambda(0)$. \\
\bottomrule\end{tabularx}
\caption{Conventions used throughout. The spatial imprint $u=v$ is supplied
in both branch constructions.}
\end{table}

\section{Geometry, spectra and Hamiltonian}\label{sec:geometry}

Let $\K=L^2(S^2,\Sigma;\dd\Omega)$ and
$\HH=L^2(\R,\dd v;\K)$. We use a nonpositive scalar Laplacian and
a self-adjoint Dirac operator. The sphere spectrum is
$\{\pm n:n\ge1\}$ with multiplicity $2n$ at each sign
\cite[(3.40), (3.54)]{CamporesiHiguchi1996}.
Set $A=D^2$ and let $P_n$ be its rank-$4n$ projection.
The chirality satisfies $\Gamma^2=I$, $\Gamma^*=\Gamma$ and
$\Gamma D=-D\Gamma$.

\begin{proposition}[Geometric square and full spectral resolution]
\label{prop:spectrum}
On the complete product cylinder with metric
$\dd v^2+\tfrac14g_{S^2}$,
\begin{equation}\label{eq:curvature}
 L_0=-\partial_v^2+4A=(\nabla^C)^*\nabla^C+2.
\end{equation}
Under the unitary Fourier transform
\begin{equation}\label{eq:fourier}
 \widehat\psi(k)=(2\pi)^{-1/2}\int_\R\ee^{-\ii kv}\psi(v)\dd v,
\end{equation}
the complete signed Dirac channels are
\begin{equation}\label{eq:blocks}
 M_n(k)=\begin{pmatrix}2n&k\\k&-2n\end{pmatrix},\quad
 E_n(k)=\sqrt{k^2+4n^2},\quad
 \Pi_\pm(n,k)=\frac12\left(I\pm\frac{M_n(k)}{E_n(k)}\right).
\end{equation}
There are $2n$ copies of each signed channel at fixed $k$, and $4n$
states for the square. The maximal domains are
\begin{align}
 \Dom\Dc&=\left\{\psi\in\HH:
 \sum_n\int_\R(k^2+4n^2)\norm{P_n\widehat\psi(k)}^2\dd k<\infty\right\},
 \label{eq:Ddomain}\\
 \Dom L_0^p&=\left\{\psi\in\HH:
 \sum_n\int_\R(k^2+4n^2)^{2p}\norm{P_n\widehat\psi(k)}^2\dd k<\infty\right\},
 \label{eq:Ldomain}
\end{align}
for every fixed integer $p\ge0$.
\end{proposition}
\begin{proof}
In a unit-sphere orthonormal frame $e_a$, the cylinder frame is
$\partial_v,2e_a$. The product connection has no mixed radial terms, so
its connection Laplacian is
$-\partial_v^2+4\nabla_{S^2}^*\nabla_{S^2}$.
The cross term in $\Dc^2$ vanishes by anticommutation.
The curvature formula for the self-adjoint convention is
$D^2=\nabla_{S^2}^*\nabla_{S^2}+\tfrac12$; equivalently the cylinder
has scalar curvature $8$ and contributes $8/4=2$
\cite[(5.35)--(5.36)]{CamporesiHiguchi1996}.
Choose $D\chi^+_{n,a}=n\chi^+_{n,a}$ and
$\chi^-_{n,a}=\Gamma\chi^+_{n,a}$, $1\le a\le2n$.
Fourier transformation gives \eqref{eq:blocks}, whose square is
$(k^2+4n^2)I$. Real Hermitian multiplication on the maximal weighted
domain is self-adjoint, giving \eqref{eq:Ddomain}--\eqref{eq:Ldomain}.
Complete-shell cutoffs converge in graph norm. For finitely many shells
the graph norms are finite sums of radial Sobolev norms; smooth cutoff
and approximation then give the core
$C_c^\infty(\R)\otimes\K_{\rm fin}$.
\end{proof}

These are continuous radial spectra:
\[
 \sigma(\Dc)=(-\infty,-2]\cup[2,\infty),\qquad
 \sigma(L_0)=[4,\infty).
\]
Each channel dispersion is monotone on either open half-line of $k$.
The thresholds have preimages of measure zero; the spectral measures
are absolutely continuous, including their integrable threshold
singularities. Plane and cosine modes are generalized modes, not
normalizable radial states.

The metric volume is $\dd v\dd\Omega/4$. Multiplication of product-measure
spinors by $2$ is its unitary identification with our Hilbert convention.
With physical length unit $\ell_0>0$ and mass $m_{\rm phys}>0$, the
declared kinetic Hamiltonian is
\[
 H_{\rm kin}=\varepsilon L_0,\qquad
 \varepsilon=\frac{\hbar^2}{2m_{\rm phys}\ell_0^2}.
\]
The curvature contribution $2\varepsilon$ is already in this expression.
Its unitary evolution is $\ee^{-\ii\varepsilon tL_0/\hbar}$.
The distinct heat evolution has norm
$\norm{\ee^{-\tau L_0}}=\ee^{-4\tau}$ and is smoothing in left energy
for $\tau>0$. Its integral is $L_0^{-1}$ with norm $1/4$.
There is no normalized full-cylinder Gibbs state: the positive diagonal
heat-trace density is constant over an infinite radial line.

For comparison, the scalar operator on the same metric is
\begin{equation}\label{eq:scalaroperator}
 \Lsc=1-\Delta_C=1-\partial_v^2-4\Delta_{S^2}.
\end{equation}
Its degree-$\ell$ channel has multiplier $k^2+(2\ell+1)^2$ and
multiplicity $2\ell+1$, including $\ell=0$. Its maximal power domains
use the corresponding weight $[k^2+(2\ell+1)^2]^{2p}$.
Thus the scalar and spinor channel decompositions are different even
though their metrics agree.

On the even sector the half-line cosine transform is
$\mathcal C h(k)=\sqrt{2/\pi}\int_0^\infty h(v)\cos(kv)\dd v$.
The even extension $h\mapsto h(|v|)/\sqrt2$ is unitary. In particular
an even full-line $f$ satisfies
$\mathcal C(\sqrt2 f|_{(0,\infty)})=\sqrt2\widehat f$ for $k>0$.
Evenness reduces $L_0$ and $\Lsc$, but does not reduce $\Dc$ because
the first derivative changes parity. No wall or doubled boundary Green
kernel is introduced by using this transform.

\section{From the Green response to the theta source}\label{sec:theta}

\subsection{Multiplicity cancellation and angular energy}
The decaying full-line Green kernel in shell $n$ is
\begin{equation}\label{eq:green}
 g_n(v,w)=\frac{\ee^{-2n|v-w|}}{4n},\qquad
 g_n(v,v)=\int_\R\frac{\dd k}{2\pi(k^2+4n^2)}=\frac1{4n}.
\end{equation}
Its derivative has jump $-1$, fixing the inverse of
$-\partial_v^2+4n^2$. Consequently the angular diagonal is the bounded
operator $W=(4\sqrt A)^{-1}$; it is not trace class. Damping makes it
trace class and cancels the complete multiplicity:
\begin{equation}\label{eq:Q}
 \Tr_\K(W\ee^{-\beta A})
 =\sum_{n\ge1}4n\frac{\ee^{-\beta n^2}}{4n}=Q(\beta),\qquad\beta>0.
\end{equation}
Every fixed $\beta$ derivative converges uniformly for $\beta\ge\beta_0>0$.
This is an angular trace of a full radial Green diagonal. It is independent
of the probe location and includes all radial momenta.

For the optional filter $B_s=P_1+s(I-P_1)$, $0\le s\le1$, set
\[
 Q_s(\beta)=\Tr(B_sW\ee^{-\beta A}),\qquad
 \rho_{s,\beta}=Q_s(\beta)^{-1}B_sW\ee^{-\beta A}.
\]
These are positive angular density operators with finite energy moments,
and direct spectral summation gives
\begin{equation}\label{eq:energymoments}
 (-1)^rQ_s^{(r)}(\beta)=Q_s(\beta)\Tr(\rho_{s,\beta}A^r),\qquad r\ge0.
\end{equation}
The identities concern an angular state; they do not require a full-space
thermal state or a detector construction. The energy-observable dictionary
below is also recorded in \cite{MC11Observables}; its trace algebra is
given here without importing a pending physical-preparation claim.

\subsection{The prescribed scan and complete reflection}
Supply the scan and differential readout
\begin{equation}\label{eq:scan}
 \beta(u)=\pi\ee^{4u},\qquad p(u)=\ee^uQ(\beta(u)),\qquad
 \Phi(u)=\tfrac18(p''(u)-p(u)).
\end{equation}
Differentiating the complete series gives \eqref{eq:source}. More
generally the filtered readout is
\begin{align}
 \Phi_s(u)&=\ee^u Q_s(\beta)
 \Tr\!\left[\rho_{s,\beta}(2\beta^2A^2-3\beta A)\right],\label{eq:sourceenergy}\\
 \Phi_s'(u)&=\ee^u Q_s(\beta)
 \Tr\!\left[\rho_{s,\beta}(-8\beta^3A^3+30\beta^2A^2-15\beta A)\right].
 \label{eq:sourceenergyprime}
\end{align}
Indeed $p''-p=\ee^u(16\beta^2Q''+24\beta Q')$; one more derivative
gives the cubic polynomial. The quadratic energy observable is positive
for $u\ge0$, since $\beta n^2\ge\pi>3/2$. Positivity is not asserted
for all $\beta>0$.

For the complete unit-weight source, put
$\Theta(R)=1+2Q(\pi R^2)$. Jacobi inversion gives
$\Theta(R)=R^{-1}\Theta(1/R)$, $R>0$
\cite[20.7.32, at zero argument]{DLMFThetaModular}.
With $R=\ee^{2u}$ this is
\begin{equation}\label{eq:reflectionp}
 p(u)-p(-u)=-\sinh u.
\end{equation}
The complete series is smooth on each real compact interval. Applying
$(\partial_u^2-1)/8$ to \eqref{eq:reflectionp} proves smooth evenness
of $\Phi$. For $u\ge0$ every term in \eqref{eq:source} is strictly
positive, hence $\Phi>0$ on $\R$. The added $1$ in $\Theta$ is the
missing zero mode and is essential to this reflection calculation.
An arbitrary shell filter does not retain this identity.

The Hamiltonian produces \eqref{eq:Q}; it does not produce the
position-dependent scan \eqref{eq:scan}. In particular translating the
Green probe cannot imprint the source profile, because
\eqref{eq:green}'s diagonal is independent of $v$.

\subsection{Bounds for the complete source at every finite order}
\begin{lemma}[Derivative and tail bounds]\label{lem:tails}
Set $x=\pi\ee^{4u}$ for $u\ge0$, and define polynomials
\begin{equation}\label{eq:recurrence}
 \mathsf P_0(x)=2x^2-3x,\qquad
 \mathsf P_{j+1}(x)=(1-4x)\mathsf P_j(x)+4x\mathsf P_j'(x).
\end{equation}
Then $\deg \mathsf P_j=j+2$ and
$\Phi^{(j)}(u)=\ee^u\sum_{m\ge1}\ee^{-m^2x}\mathsf P_j(m^2x)$.
If $a_j$ is the absolute coefficient sum of $\mathsf P_j$, then
\begin{align}
 C_d(N)&=\sum_{m>N}m^{2d}\ee^{-\pi(m^2-1)}\longrightarrow0,\label{eq:Ctail}\\
 \ee^u\sum_{m>N}\ee^{-m^2x}|\mathsf P_j(m^2x)|
 &\le a_j C_{j+2}(N)\ee^u x^{j+2}\ee^{-x}.
 \label{eq:tailbound}
\end{align}
Writing $r=|u|$, $x=\pi\ee^{4r}$, $b_j=a_jC_{j+2}(0)$ and
$\kappa=2-3/\pi>0$, the complete smooth source obeys
\begin{equation}\label{eq:twosidedbounds}
 \kappa\ee^r x^2\ee^{-x}\le\Phi(u),\qquad
 |\Phi^{(j)}(u)|\le b_j\ee^r x^{j+2}\ee^{-x}.
\end{equation}
\end{lemma}
\begin{proof}
Differentiate one summand using $x'=4x$ to obtain
\eqref{eq:recurrence}. Because $x\ge\pi>1$,
$|\mathsf P_j(m^2x)|\le a_jm^{2(j+2)}x^{j+2}$, and
$\ee^{-m^2x}\le\ee^{-x}\ee^{-\pi(m^2-1)}$.
Gaussian summability proves \eqref{eq:Ctail}--\eqref{eq:tailbound}.
The first source term gives the lower bound, since
$2x^2-3x\ge\kappa x^2$.
Complete reflection transfers the bounds to the other half-line.
\end{proof}

For any fixed nonnegative integers $a,d$ and finite $M,b\ge0$,
\begin{equation}\label{eq:J}
 J_{a,d}(M,b)=\int_\R |u|^a
 \ee^{Mu^2+(b+1)|u|}x^d\ee^{-x}\dd u<\infty,
 \qquad x=\pi\ee^{4|u|}.
\end{equation}
Indeed the substitution $x=\pi\ee^{4|u|}$ gives
\[
 \frac{2}{4^{a+1}}\int_\pi^\infty
 [\log(x/\pi)]^a(x/\pi)^{(b+1)/4}x^{d-1}
 \exp\!\left(\frac M{16}\log^2(x/\pi)-x\right)\dd x.
\]
The negative linear term in $x$ dominates the logarithmic square and
every indicated power. This also proves convergence of the complete
integral tails. Multiplying \eqref{eq:tailbound} by these weights gives
vanishing integrated angular tails at each fixed order.

All constants may depend on the chosen derivative order. These estimates
do not give uniform bounds as the order tends to infinity. Also, a
reflected finite prefix of the source need not be smooth at zero. We use
such prefixes only for value or moment tail estimates; smoothness and
vanishing odd derivatives of the exact source come from complete reflection.
These are the bounds underlying the shared comparison
\cite[Section 6]{TNDParent}.
